% Part: history
% Chapter: biographies
% Section: alonzo-church

\documentclass[../../../include/open-logic-section]{subfiles}

\begin{document}

\olfileid{his}{bio}{chu}

\olsection{آلونزو چرچ}

\olphoto{church-alonzo}{Alonzo Church}

آلونزو چرچ در 14 ژوئن 1903 در واشنگتن دی‌سی به دنیا آمد. در اوایل
کودکی، حادثه‌ای با تفنگ بادی او را از یک چشم نابینا کرد. چرچ در سال
1920 مدرسهٔ مقدماتی را در کنتیکت به پایان رساند و همان سال تحصیلات
دانشگاهی خود را در پرینستون آغاز کرد. او دورهٔ دکتری را در 1927 به
پایان رساند. چرچ پس از دو سال اقامت در خارج به پرینستون بازگشت. او به
ادب و دقت بسیار شهرت داشت. نوشته‌هایش روی تخته بی‌نقص بود و مقاله‌های
مهم را با پوشاندن دقیقشان با سیمان دوکو، نوعی چسب شفاف، نگهداری می‌کرد.
بیرون از فعالیت‌های دانشگاهی از خواندن مجله‌های علمی‌تخیلی لذت می‌برد
و اگر نادرستی‌ای در نوشته می‌یافت، از نامه‌نوشتن به سردبیران ابایی نداشت.

دستاوردهای دانشگاهی چرچ بزرگ بودند. او همراه با دانشجویانش استیون
کلینی و بارکلی راسر، مستقل از توسعهٔ ماشین تورینگ به دست آلن تورینگ،
نظریه‌ای دربارهٔ محاسبه‌پذیری مؤثر، یعنی حساب لامبدا، پروراند. این دو
تعریف محاسبه‌پذیری هم‌ارزند و به آنچه اکنون \emph{تز چرچ--تورینگ}
نامیده می‌شود می‌انجامند: تابعی بر اعداد طبیعی به‌طور مؤثر محاسبه‌پذیر
است اگر و تنها اگر با ماشین تورینگ یا حساب لامبدا محاسبه‌پذیر باشد.
او همچنین چیزی را اثبات کرد که اکنون \emph{قضیهٔ چرچ} نام دارد: مسئلهٔ
تصمیم برای اعتبار فرمول‌های مرتبهٔ اول حل‌ناپذیر است.

چرچ تا سالخوردگی به کار خود ادامه داد. در سال 1967 پرینستون را به مقصد
دانشگاه کالیفرنیا در لس‌آنجلس ترک کرد و تا بازنشستگی در 1990 استاد آنجا
بود. چرچ در 1 اوت 1995 در 92سالگی درگذشت.

\begin{reading}
برای زندگی‌نامه‌ای کوتاه از چرچ، \citet{EndertonND} را ببینید. نوشته‌های
اصلی چرچ دربارهٔ حساب لامبدا و مسئلهٔ تصمیم، یا تز چرچ، در
\citet{Church1936,Church1936a} آمده‌اند. \citet{Aspray1984} گفت‌وگویی
با چرچ دربارهٔ جامعهٔ ریاضی پرینستون در دههٔ 1930 را ثبت کرده است.
چرچ از سال 1936 تا 1979 مجموعه‌ای از نقدهای کتاب را برای
\emph{Journal of Symbolic Logic} نوشت. همهٔ آن‌ها در وبگاه جان
مک‌فارلین بایگانی شده‌اند \citep{MacFarlane2015}.
\end{reading}
\end{document}
